% Part: computability
% Chapter: recursive-functions
% Section: partial-functions

\documentclass[../../../include/open-logic-section]{subfiles}

\begin{document}

\olfileid[hi]{cmp}{rec}{par}
\olsection{आंशिक पुनरावर्ती फलन}


पुनरावर्ती फलनों की परिभाषा को प्रेरित करने के लिए ध्यान दें कि हमारा यह
प्रमाण कि ऐसे संगणनीय फलन हैं जो आदिम पुनरावर्ती नहीं हैं, वास्तव में
कहीं अधिक स्थापित करता है। तर्क सरल था: हमने केवल इस तथ्य का उपयोग किया
कि फलनों $f_0,f_1,\dots$ को इस प्रकार परिगणित करना संभव है कि $x$ और
$y$ के फलन के रूप में $f_x(y)$ संगणनीय हो। अतः तर्क \emph{फलनों के हर
ऐसे वर्ग पर लागू होता है जिसे इस ढंग से परिगणित किया जा सके}। इससे हम
दुविधा में पड़ते हैं: हम संगणनीय फलनों का स्पष्ट वर्णन करना चाहते हैं,
किंतु संगणनीय फलनों के किसी संग्रह का कोई भी स्पष्ट वर्णन सर्वसमावेशी
नहीं हो सकता!

इससे निकलने का मार्ग \emph{आंशिक} फलनों को भूमिका देना है। हम देखेंगे कि
आंशिक संगणनीय फलनों को परिगणित करना \emph{संभव है}।\iftag{TMs}{ वास्तव
में हम लगभग पहले से जानते हैं कि ऐसा है, क्योंकि ट्यूरिंग मशीनों को
व्यवस्थित ढंग से परिगणित करना संभव है।}{} आगे हम अपने विकर्ण तर्क पर
लौटेंगे और देखेंगे कि आंशिक फलनों को सम्मिलित करने पर वह क्यों नहीं चलता।

अब प्रश्न यह है: सभी आंशिक पुनरावर्ती फलन पाने के लिए हमें आदिम
पुनरावर्ती फलनों में क्या जोड़ना होगा? हमें दो काम करने हैं:
\begin{enumerate}
\item आदिम पुनरावर्ती फलनों की अपनी परिभाषा को इस प्रकार बदलना कि उसमें
  आंशिक फलनों की भी अनुमति हो।
\item परिभाषा में कुछ \emph{जोड़ना}, ताकि उसमें कुछ नए आंशिक फलन आएँ।
\end{enumerate}

पहला काम सरल है। पहले की तरह हम शून्य, उत्तराधिकारी और प्रक्षेपणों से
आरंभ करेंगे और संयोजन तथा आदिम पुनरावर्तन के अधीन संवृत करेंगे। केवल यह
अंतर है कि संयोजन और आदिम पुनरावर्तन की परिभाषाओं को इस संभावना की अनुमति
देने के लिए बदलना होगा कि परिभाषा के कुछ पद परिभाषित न हों। यदि $f$ और
$g$ आंशिक फलन हैं, तो $f(x) \fdefined$ का अर्थ होगा कि $f$, $x$ पर
परिभाषित है, अर्थात् $x$, $f$ के प्रांत में है; और $f(x)
\fundefined$ का अर्थ इसका उलटा होगा, अर्थात् $f$,~$x$ पर परिभाषित नहीं
है। हम $f(x) \simeq g(x)$ का अर्थ लेंगे कि या तो $f(x)$ और $g(x)$ दोनों
अपरिभाषित हैं, अथवा दोनों परिभाषित और बराबर हैं। अधिक जटिल पदों के लिए भी
हम इन संकेतों का उपयोग करेंगे। हम यह परिपाटी अपनाएँगे कि यदि $h$ और
$g_0$, \dots,~$g_k$ सभी आंशिक फलन हों, तो
\[
h(g_0(\vec x),\dots,g_k(\vec x))
\]
तभी और केवल तभी परिभाषित है जब प्रत्येक $g_i$, $\vec x$ पर परिभाषित हो
और $h$, $g_0(\vec x)$, \dots,~$g_k(\vec x)$ पर परिभाषित हो। इस समझ के
साथ आंशिक फलनों के लिए संयोजन और आदिम पुनरावर्तन की परिभाषाएँ ठीक ऊपर
जैसी हैं, सिवाय इसके कि हमें ``$=$'' को ``$\simeq$'' से बदलना होगा।

आंशिक फलन पाने के लिए आदिम पुनरावर्ती फलनों की परिभाषा में हम
\emph{अपरिबद्ध खोज संकारक} जोड़ेंगे। यदि $f(x,\vec z)$ प्राकृतिक
संख्याओं पर कोई आंशिक फलन है, तो $\mu x \; f(x,\vec z)$ को यह मान
परिभाषित करें:
\begin{quote}
  वह लघुतम $x$ जिसके लिए $f(0,\vec z), f(1,\vec z), \dots, f(x,\vec
  z)$ सभी परिभाषित हैं और $f(x,\vec z) = 0$ है, यदि ऐसा कोई $x$
  अस्तित्व में हो।
\end{quote}
इस समझ के साथ कि अन्यथा $\mu x \; f(x,\vec z)$ अपरिभाषित है। इससे
$\mu x \; f(x,\vec z)$ अद्वितीय रूप से परिभाषित होता है।

\begin{explain}
ध्यान दें कि हमारी परिभाषा में\iftag{TMs}{ ट्यूरिंग मशीनों,}{} कलनविधियों
या किसी विशिष्ट संगणकीय प्रतिरूप का कोई संदर्भ नहीं है। किंतु संयोजन और
आदिम पुनरावर्तन की तरह अपरिबद्ध खोज के पीछे भी एक संक्रियात्मक, संगणकीय
सहज-बोध है। किसी आंशिक फलन की संगणनीयता के संदर्भ में, जिन तर्कों पर फलन
अपरिभाषित है वे उन निवेशों के संगत हैं जिनके लिए संगणन रुकता नहीं।
$\mu x \;
f(x,\vec z)$ को संगणित करने की प्रक्रिया यह होगी:
$f(0,\vec z), f(1,\vec z),
f(2,\vec z)$ को तब तक संगणित करें जब तक मान
0 न लौटे। किंतु यदि बीच का कोई संगणन न रुके, तो
$\mu x \; f(x,\vec z)$ का संगणन भी नहीं रुकता।
\end{explain}

यदि $R(x,\vec z)$ कोई संबंध है, तो $\mu x \; R(x,\vec z)$ को
$\mu x \; (1 \tsub \Char{R}(x,\vec z))$ परिभाषित किया जाता है। दूसरे
शब्दों में, $\mu x \;
R(x,\vec z)$ उस लघुतम $x$ को लौटाता है जिसके लिए
$R(x,\vec z)$ मान्य है। अतः यदि $f(x,\vec z)$ कोई पूर्ण फलन है, तो
$\mu x \; f(x,\vec
z)$ वही है जो $\mu x \; (f(x,\vec z) = 0)$। किंतु
ध्यान दें कि हमारी मूल परिभाषा अधिक सामान्य है, क्योंकि वह इस संभावना
की अनुमति देती है कि $f(x,\vec z)$ सर्वत्र परिभाषित न हो (इसके विपरीत,
किसी संबंध का अभिलाक्षणिक फलन सदा पूर्ण होता है)।

\begin{defn}
\emph{आंशिक पुनरावर्ती फलनों} का समुच्चय प्राकृतिक संख्याओं से प्राकृतिक
संख्याओं में (विभिन्न स्थान-संख्याओं वाले) आंशिक फलनों का वह लघुतम
समुच्चय है जिसमें शून्य, उत्तराधिकारी और प्रक्षेपण आते हैं तथा जो संयोजन,
आदिम पुनरावर्तन और अपरिबद्ध खोज के अधीन संवृत है।
\end{defn}

निस्संदेह, कुछ आंशिक पुनरावर्ती फलन पूर्ण होंगे, अर्थात् प्रत्येक तर्क
के लिए परिभाषित होंगे।

\begin{defn}
\ollabel{defn:recursive-fn}
\emph{पुनरावर्ती फलनों} का समुच्चय उन आंशिक पुनरावर्ती फलनों का समुच्चय
है जो पूर्ण हैं।
\end{defn}

किसी पुनरावर्ती फलन को कभी-कभी ``पूर्ण पुनरावर्ती'' कहा जाता है, ताकि इस
बात पर बल दिया जा सके कि वह सर्वत्र परिभाषित है।

\end{document}

